\newpage
\keyframe{Document Information and Key}

This document is shared with Malcolm Joyce (project supervisor), Fleur Hughes (programme manager), James Lewis (programme coordinator), Dan Doherty (taught programme lead) and Alyanna Dela Cruz (PLANET PGR at Lancaster). 

It is not a formal document for submission and the content is not necessarily original. For example, the majority of chapter \ref{nucdat}: \textit{Nuclear Data} can be taken to be skeleton notes from Jenkins' and Wood's \textit{Nuclear Data} series. It is to be used as a personal reference tool, diary and as proof of engagement. As such, formal references will not be given, only in the case of copied or adapted diagrams and where it is useful for self-reference. Instead of references, an informal list of sources will be given in the introduction to each chapter. 

AI will not be used for generating or proof-reading text, writing any LaTeX markup or for formatting. For time's sake, it will be used for reproducing editable diagrams where I would like to adapt them in some way. For example, figure \ref{rms_vs_A} was reproduced by using Claude Sonnet 5.5. Where this is done, a \anth symbol will be shown (closest dingbat to Anthropic's logo). I have chosen to work with Claude as it is familiar. I have searched for Lancaster's recommended AI tool, but cannot find any definitive answer. In 2023, Lancaster's \href{https://wp.lancs.ac.uk/lu-digitaldigest/2023/12/01/ai-tools-a-christmassy-showcase/}{\ul{digital digest}} included mention of an online training course, \href{https://12daysofai.myblog.arts.ac.uk/daily-activities/}{\ul{12 Days of AI}}, Claude is recommended on Day 4.

\begin{example}{1}{Worked Example}
\textbf{Solutions to assigned exercises. The `difficulty dots' match the asterisks given in \textit{Primer}. For multiple parts of varying difficulties, the highest number is taken.}
\end{example}


\begin{derivation}{2}{Derivation}
\textbf{\raggedright Derivations of equations.}
\end{derivation}

\begin{historical}{Historical Context}
\textbf{\raggedright Historical context and background information.}
\end{historical}




\vfill